\section{Learning roadmap}

These notes follow two sources supplied by the research supervisor:
the thesis \emph{Solitons in Bose--Einstein Condensates} by Matthias S\"ohn,
and the article \emph{Solitons in Bose--Einstein condensates} by
R. Balakrishnan and I. I. Satija.

Each lesson has one conceptual job, one central derivation, and one checkpoint.

\subsection{Block A: From ordinary quantum mechanics to a condensate}
\begin{enumerate}[label=A\arabic*.]
  \item Identical bosons and symmetric two-particle states.
  \item Occupation numbers and Fock states.
  \item Creation and annihilation operators.
  \item Field operators and the density operator.
  \item Ideal Bose gas and the Bose--Einstein distribution.
  \item Saturation of excited states and the critical temperature.
  \item Condensation, one-body density matrix, and coherence.
  \item Global $U(1)$ symmetry and the condensate order parameter.
\end{enumerate}

\subsection{Block B: Interactions and the Gross--Pitaevskii equation}
\begin{enumerate}[label=B\arabic*.]
  \item First- and second-quantized many-boson Hamiltonians.
  \item Low-energy $s$-wave scattering and scattering length.
  \item Contact pseudopotential and coupling $g$.
  \item Exact Heisenberg equation for the Bose field.
  \item Mean-field approximation and derivation of the GPE.
  \item Energy functional, normalization, and chemical potential.
\end{enumerate}

\subsection{Block C: Preparing the soliton problem}
\begin{enumerate}[label=C\arabic*.]
  \item Uniform condensate and stationary solutions.
  \item Density--phase representation and superfluid velocity.
  \item Linearization and Bogoliubov excitations.
  \item Sound speed and healing length.
  \item Reduction from three dimensions to a quasi-one-dimensional GPE.
  \item Dimensionless nonlinear Schr\"odinger equation.
\end{enumerate}

\subsection{Block D: Gross--Pitaevskii solitons}
\begin{enumerate}[label=D\arabic*.]
  \item Dispersion--nonlinearity balance and travelling-wave reduction.
  \item Bright soliton for attractive interactions.
  \item Dark soliton for repulsive interactions.
  \item Grey soliton: velocity, depth, and phase jump.
  \item Energy, momentum, stability, and numerical propagation.
\end{enumerate}

\subsection{Block E: Beyond the ordinary GPE}
\begin{enumerate}[label=E\arabic*.]
  \item Hartree many-body description of a bright soliton.
  \item Attractive delta gas and the Bethe ansatz.
  \item Localized quantum soliton states and centre-of-mass spreading.
  \item Correlations and the limits of mean field.
  \item Hard-core bosons and the Bose--Hubbard model.
  \item Spin-$\tfrac12$ mapping and spin-coherent states.
  \item Dark, antidark, and persistent hard-core-boson solitons.
  \item Particle--hole symmetry and energy--momentum dispersion.
\end{enumerate}

The first working lesson is A1. It introduces only symmetrization and uses a
two-boson example; Fock-space notation begins in A2.
